\documentclass[12pt]{article}
\usepackage[utf8]{inputenc}
\usepackage[spanish]{babel}
\usepackage{amsmath}
\usepackage{amsfonts}
\usepackage{amssymb}
\usepackage{geometry}
\geometry{letterpaper, margin=2.5cm}

\begin{document}

\begin{center}
    \Large \textbf{Evaluación de Examen 1: Unidad I} \\
    \large Física para Ciencias de la Salud (005-1713) \\
    \vspace{0.5cm}
    \textbf{Estudiante:} Mariana Salazar \quad \textbf{C.I.:} 33.570.679
    \vspace{0.3cm} \\
    \large \textbf{Calificación Final: 8/10 puntos}
\end{center}

\vspace{0.5cm}

\section*{Análisis de Resultados}

\subsection*{1. Ángulo plano (Conversión a radianes)}
\textbf{Procedimiento y resultado:} La estudiante utiliza una regla de tres basándose en la equivalencia $90^\circ = \frac{\pi}{2}$ rad. La lógica matemática aplicada es correcta y la lleva al resultado válido de $0,41\pi$ rad (el cual es una aproximación decimal de $\frac{5\pi}{12}$). \\
\textbf{Veredicto:} \textbf{Correcto} (2/2 pts).

\subsection*{2. Sistemas de coordenadas (Longitud del hueso)}
\textbf{Procedimiento y resultado:} Buena resolución gráfica y analítica. Dibuja el triángulo rectángulo, identifica las medidas de los catetos ($a=6$, $b=8$) y aplica el Teorema de Pitágoras. Existe un ligero error de notación al final del desarrollo (escribe $c^2 = \sqrt{100}$ y remata con $c^2 = 10$, cuando debió escribir $c=10$), pero las operaciones matemáticas y el resultado final son correctos. \\
\textbf{Veredicto:} \textbf{Correcto} (2/2 pts).

\subsection*{3. Vectores (Fuerza resultante)}
\textbf{Procedimiento y resultado:} Entiende el concepto de suma vectorial y agrupa correctamente las componentes $\hat{i}$ y $\hat{j}$. Sin embargo, comete un \textbf{error de signo} evidente en la componente $\hat{j}$: en lugar de sumar $(25 + 35)\hat{j}$, realiza una resta $(25 - 35)\hat{j}$, lo que la lleva a un vector incorrecto de $30\hat{i} - 10\hat{j}$ N. \\
\textbf{Veredicto:} \textbf{Parcialmente correcto} (Planteamiento válido, error de signo al operar) (1/2 pts).

\subsection*{4. Potencias de diez (Notación científica y prefijo)}
\textbf{Procedimiento y resultado:} Expresa la medida decimal adecuadamente empleando una potencia de base diez ($12,5 \times 10^{-6}$). Sin embargo, la estudiante no completó la segunda parte de la instrucción: omitió identificar e indicar explícitamente el prefijo del Sistema Internacional correspondiente (micrómetros). \\
\textbf{Veredicto:} \textbf{Incompleto / Parcialmente correcto} (1/2 pts).

\subsection*{5. Sistemas de unidades (Conversión de densidad)}
\textbf{Procedimiento y resultado:} Aunque el planteamiento inicial de los factores de conversión resulta un poco confuso, la estudiante desglosa acertadamente el procedimiento aritmético en dos pasos lógicos: divide primero entre 1000 (para transformar los gramos a kilogramos) y posteriormente multiplica por $1.000.000$ (para llevar los $\text{cm}^3$ a $\text{m}^3$). Obtiene el resultado numérico exacto de $1030 \text{ kg/m}^3$. \\
\textbf{Veredicto:} \textbf{Correcto} (2/2 pts).

\vspace{0.5cm}
\section*{Comentarios Generales y Sugerencias}
La estudiante demuestra un buen dominio en el cálculo aritmético básico y la geometría aplicada a coordenadas. Se recomiendan los siguientes puntos de mejora:
\begin{itemize}
    \item \textbf{Cuidado con los signos:} Es fundamental repasar la suma algebraica de vectores para evitar restar los componentes cuando el enunciado solicita la fuerza resultante (Pregunta 3).
    \item \textbf{Notación formal:} Se aconseja mantener rigor en la notación matemática, como no arrastrar el exponente cuadrático ($c^2$) después de haber aplicado la raíz cuadrada (Pregunta 2).
\end{itemize}

\end{document}